\chapter{Диффузия}\label{ch:diffusion}

\section{Прыжки атомов и случайное блуждание}

Атом в кристалле большую часть времени колеблется около своего места и лишь изредка, когда тепловая флуктуация
даёт ему энергию $\Delta G_m$, перепрыгивает в соседнюю позицию. Частота прыжков
\begin{equation}
\Gamma = \nu\, e^{-\Delta G_m/k_BT},
\end{equation}
где $\nu\sim10^{13}$~с$^{-1}$ --- частота колебаний. Водород в цирконии --- атом внедрения: он сидит в
тетраэдрических междоузлиях и прыгает между ними, не нуждаясь в вакансиях. Поэтому он очень подвижен.

Каждый прыжок --- в случайную сторону. После $n$ прыжков длиной $a$ средний сдвиг равен нулю, а средний
\emph{квадрат} сдвига растёт линейно (рис.~\ref{fig:f3_walk}):
\begin{equation}
\langle x^2\rangle = n a^2 \quad(\text{одномерно}),\qquad n = \Gamma t .
\end{equation}
Определим коэффициент диффузии так, чтобы $\langle x^2\rangle = 2Dt$ в одном измерении ($6Dt$ в трёх). Тогда
\begin{keyf}[Коэффициент диффузии]
\[
D = \frac{\Gamma_{\text{пол}} a^2}{6} = D_0\,e^{-Q/RT},
\]
где $\Gamma_{\text{пол}}$ --- полная частота прыжков во всех направлениях, $Q$ --- энергия активации миграции.
\end{keyf}

\fig{f3_walk}{а) Траектории пяти атомов по 400 случайных прыжков: каждая непредсказуема, но типичный уход ---
порядка $\sqrt n\,a$ (пунктирная окружность). б) Распределение смещений 20\,000 одномерных блужданий после $n$
прыжков: гистограммы ложатся на гауссианы с $\langle x^2\rangle = na^2$.}

\begin{idea}
Главное следствие: расстояние растёт как \emph{корень} из времени. Чтобы уйти вдвое дальше, нужно вчетверо
больше времени. Диффузионная длина $\sqrt{Dt}$ --- универсальная мерка «куда успеет дойти атом за время $t$».
\end{idea}

\section{Законы Фика}

На макроуровне случайные прыжки дают поток из мест, где атомов больше, туда, где их меньше (первый закон Фика):
\begin{equation}
\mathbf J = -D\,\nabla c .
\end{equation}
Сохранение числа атомов $\partial c/\partial t = -\nabla\cdot\mathbf J$ даёт второй закон:
\begin{keyf}[Уравнение диффузии]
\[
\frac{\partial c}{\partial t} = \nabla\cdot(D\nabla c) \;=\; D\,\nabla^2 c\quad(D = \text{const}).
\]
\end{keyf}

Два решения, которые стоит знать наизусть (рис.~\ref{fig:f3_profiles}):
\begin{itemize}
\item \textbf{тонкий слой} количества $M$ в точке $x=0$ расплывается гауссианой
  \[c(x,t) = \frac{M}{\sqrt{4\pi Dt}}\,e^{-x^2/4Dt};\]
\item \textbf{постоянная концентрация $c_s$ на поверхности} полубесконечного тела:
  \[c(x,t) = c_s\,\operatorname{erfc}\Bigl(\frac{x}{2\sqrt{Dt}}\Bigr).\]
  На глубине $x = 2\sqrt{Dt}$ концентрация уже
  $0.16\,c_s$.
\end{itemize}

\fig{f3_profiles}{а) Тонкий слой расплывается гауссианой. б) Профиль от поверхности с постоянной концентрацией
--- функция erfc; точки --- глубина $2\sqrt{Dt}$, где концентрация 16\,\% от поверхностной. Во всех решениях
масштаб длины --- $\sqrt{Dt}$.}

\section{Водород в цирконии: насколько он подвижен}

Для водорода в $\alpha$-цирконии (Kearns, 1972)
\[
D \approx 7.9\cdot10^{-7}\exp\Bigl(-\frac{44.4~\text{кДж/моль}}{RT}\Bigr)~\text{м}^2/\text{с}.
\]
При 300\degC{} $D\approx 7\cdot10^{-11}$~м$^2$/с, при 400\degC{} --- $3\cdot10^{-10}$. За минуту водород
проходит около 0.1~мм, за час --- 0.5--1~мм, то есть всю стенку трубы (рис.~\ref{fig:f3_dh}).

\begin{figure}[htbp]\centering
\includegraphics[width=0.66\linewidth]{f3_dh}
\caption{Диффузионная длина водорода в $\alpha$-Zr $\sqrt{Dt}$ от температуры для 1~с, 1~мин, 1~ч и 1~сут против
характерных размеров структуры. При охлаждении 0.5--1\degC/мин водород за время прохода каждого градуса успевает
уйти на десятки микрометров --- больше расстояния между гидридами.}
\label{fig:f3_dh}
\end{figure}

\begin{zrbox}
Отсюда важный вывод для модели. При обычных скоростях охлаждения водород успевает выравнивать концентрацию на
масштабе зёрен и между соседними гидридами. Выпадение ограничено не подводом водорода издалека, а тем, где и
когда удаётся \emph{зародиться} пластинке. Зато на масштабе ниже сотни микрометров диффузия решает, какой из
соседних зародышей «перехватит» водород и вырастет.
\end{zrbox}

\section{Движущая сила --- градиент химического потенциала}

Закон Фика в форме $\mathbf J = -D\nabla c$ --- частный случай. Атомы на самом деле текут вниз по
\emph{химическому потенциалу}:
\begin{equation}
\mathbf J = -\frac{D\,c}{RT}\,\nabla\mu .
\end{equation}
Для разбавленного раствора $\mu = \mu^\infty + RT\ln c$, и $\nabla\mu = RT\nabla c/c$ --- возвращаемся к Фику. Но
если на $\mu$ влияет что-то ещё, поток может идти даже \emph{против} градиента концентрации («восходящая
диффузия»).

Самый важный для нас пример --- \textbf{напряжения}. Атом водорода расширяет решётку на свой парциальный мольный
объём $\bar V_{\mathrm H}\approx 1.7$~см$^3$/моль. В растянутой области (гидростатическое напряжение
$\sigma_h = (\sigma_{11}+\sigma_{22}+\sigma_{33})/3 > 0$) поместить его легче: работа $-\bar V_{\mathrm H}\sigma_h$
понижает его потенциал:
\begin{equation}
\mu = \mu^\infty + RT\ln c - \bar V_{\mathrm H}\,\sigma_h .
\end{equation}

\begin{keyf}[Диффузия в поле напряжений]
\[
\mathbf J = -D\Bigl(\nabla c - \frac{c\,\bar V_{\mathrm H}}{RT}\nabla\sigma_h\Bigr),
\qquad\text{в равновесии}\quad c = c_0\exp\Bigl(\frac{\bar V_{\mathrm H}\sigma_h}{RT}\Bigr).
\]
\end{keyf}

Перетекание водорода в растянутые места называют \textbf{эффектом Горского}. Он тем сильнее, чем ниже
температура и больше напряжение (рис.~\ref{fig:f3_stress}). При 300\degC{} и $\sigma_h = 300$~МПа концентрация
выше на 11\,\%. У вершины трещины, где $\sigma_h$ доходит до 2--3 пределов текучести, --- в полтора раза.
На этом стоит механизм замедленного гидридного растрескивания: водород стекается к вершине трещины, там
выпадает гидрид, гидрид трескается, трещина продвигается.

\fig{f3_stress}{а) Равновесное обогащение водородом перед вершиной трещины ($K = 10$~МПа$\sqrt{\text{м}}$,
плоская деформация). Вблизи вершины напряжение ограничено пластичностью ($\sigma_h\approx3\sigma_y$), поэтому
обогащение выходит на полку. б) Обогащение $c/c_0$ от гидростатического напряжения при трёх температурах.}

\begin{zrbox}
Внешнее окружное напряжение 100~МПа в трубе само по себе меняет концентрацию всего на 1--2\,\%: водород
перераспределяется слабо. Важнее \emph{локальные} напряжения у торцов уже выпавших гидридов --- там $\sigma_h$
доходит до сотен МПа (глава~\ref{ch:elastic}), и водород стекается к торцам. Так поле напряжений и диффузия вместе
выстраивают гидриды в цепочки.
\end{zrbox}

Есть и \textbf{термодиффузия} (эффект Соре): в градиенте температуры водород уходит в холодную сторону. В твэле
внешняя поверхность оболочки холоднее, и водород копится у неё --- так образуется гидридный обод. Формально это
ещё один вклад в $\mu$, пропорциональный теплоте переноса $Q^*$.

\section{Как решать уравнение диффузии на компьютере}

Простейшая \textbf{явная схема} на сетке с шагом $\Delta x$:
\begin{equation}
c_i^{n+1} = c_i^n + r\,\bigl(c_{i+1}^n - 2c_i^n + c_{i-1}^n\bigr),\qquad r = \frac{D\,\Delta t}{\Delta x^2}.
\end{equation}
Она устойчива, только если $r \le 1/2$ в одномерном случае ($1/4$ на квадратной сетке в 2D). При большем шаге
самая короткая гармоника (пила с периодом $2\Delta x$) умножается за шаг на $|1-4r| > 1$ и растёт без предела
(рис.~\ref{fig:f3_numerics}).

\fig{f3_numerics}{Явная схема для уравнения диффузии после 80 шагов. а) $r = 0.45$: совпадает с точным решением.
б) $r = 0.51$: чуть выше предела, и из ошибок округления уже выросла пила, которая дальше растёт
экспоненциально.}

\begin{idea}
Почему именно $1/2$? Явная схема переносит за шаг долю $r$ вещества в каждого соседа. Если $r > 1/2$, клетка
отдаёт больше, чем имела, --- и концентрация «перескакивает» через равновесие, раскачиваясь всё сильнее.
\end{idea}

Модель гидридов решает диффузию \textbf{спектральным методом}: поле раскладывается по гармоникам Фурье
$e^{i\mathbf k\cdot\mathbf x}$, и каждая затухает точно как $e^{-Dk^2\Delta t}$. Для постоянного $D$ на
периодической сетке это точное решение при любом шаге по времени. Тот же приём --- и для упругих полей
(глава~\ref{ch:elastic}).

\begin{exercise}
Оцените время, за которое водород при 350\degC{} выравнивает концентрацию по толщине стенки 0.65~мм.
А по расстоянию между гидридами 50~мкм?
\end{exercise}

\begin{exercise}
Выведите равновесное распределение $c\propto e^{\bar V\sigma_h/RT}$, приравняв поток нулю. Оцените обогащение при
$\sigma_h = 600$~МПа (у торца гидрида) при 300\degC.
\end{exercise}

\begin{exercise}
Подставьте в явную схему пробное решение $c_i^n = g^n e^{ik i\Delta x}$ и найдите множитель $g(k)$. Докажите, что
$|g|\le1$ для всех $k$ при $r\le1/2$.
\end{exercise}
